%% kapa: exam=1 round=37 year=2026 period=1 form=A subject=civil_law subject_ko=민법 seq=9 number=9
%% source: Q-Net 감정평가사 기출문제 (cst00309) artlSeq=5258924 file=raw/1008/5258924/2026년 제37회 감정평가사 1차 1교시.pdf
%% conversion: kordoc 4.24.51 -> tools/gen1.py (auto); review=proofread-sonnet; flags=none
\begin{question}{9}
불공정한 법률행위에 관한 설명으로 옳은 것은? (다툼이 있으면 판례에 따름)
\begin{choices}
  \item 경매에도 불공정한 법률행위에 관한 민법 제104조가 적용된다.
  \item 불공정한 법률행위에는 무효행위의 전환에 관한 법리가 적용될 수 있다.
  \item 불공정한 법률행위의 무효는 선의의 제3자에게 대항할 수 없다.
  \item 대리인에 의한 법률행위에서 경솔과 무경험은 본인을 기준으로 판단한다.
  \item 불공정한 법률행위로 무효인 법률행위는 추인하면 그때부터 유효로 된다.
\end{choices}
\end{question}
\answer{2}
